\subsection{相伴子群}

设 \(\mathbf{G}\) 是 \(\mathbf{R}^n\) 的任意子群（闭或不闭）。考虑集合 \(\mathbf{G}^*\) （由满足如下条件的点 \(\boldsymbol{u} = (u_i) \in \mathbf{R}^n\) 组成）：对一切 \(\boldsymbol{x} = (x_i) \in \mathbf{G}\) ，数 \((\boldsymbol{u}| \boldsymbol{x}) = \sum_{i=1}^{n} u_i x_i\) 是整数。立即可以看出， \(\mathbf{G}^*\) 是 \(\mathbf{R}^n\) 的一个子群；它称为与 \(\mathbf{G}(*)\) 相伴的子群。若 \(\mathbf{G}\) 与 \(\mathbf{H}\) 是 \(\mathbf{R}^n\) 的两个子群，且 \(\mathbf{H} \subset \mathbf{G}\) ，则显然有 \(\mathbf{G}^* \subset \mathbf{H}^*\) 。

命题 4. 子群 \(\mathbf{G}^*\) ——它与一个子群 \(\mathbf{G}\) （ \(\mathbf{R}^n\) 的）相伴——是闭的，且我们有 \((\overline{\mathbf{G}})^* = \mathbf{G}^*\) 。

对每个 \(x \in G\) ，用 \(f_x(u)\) 表示 \((u|x)\) ； \(f_x\) 是线性形式，因此是连续的。由于 \(G^*\) 是诸集合 \(f_x^{-1}(Z)\) （当 \(x\) 取遍 \(G\) 时）的交，且这些集合中的每一个都是闭的，由此得出 \(G^*\) 是闭的。另一方面，若 \(u \in G^*\) ，则有 \((u|x) \in Z\) （对一切 \(x \in G\) ），因而，由于 \(Z\) 在 \(R\) 中是闭的，有 \((u|y) \in Z\) （对一切 \(y \in \overline{G}\) ）；于是 \(u \in (G)^*\) 。但我们有 \((\overline{G})^* \subset G^*\) （因为 \(G \subset \overline{G}\) ），所以 \((\overline{G})^* = G^*\) 。

考虑 \(\mathbf{G}^*\) 在 \(\mathbf{G}\) 为闭时的结构。由第 2 小节定理 2 的推论 1，存在一个基 \((\pmb{a}_i)_{\mathbf{1} \leqslant i \leqslant n}\) （ \(\mathbf{R}^n\) 的），使得 \(\mathbf{G}\) 恰为由点

\[
\boldsymbol {x} = \sum_ {i = 1} ^ {p} t _ {i} \boldsymbol {a} _ {i} + \sum_ {j = p + 1} ^ {p + q} n _ {j} \boldsymbol {a} _ {j},
\]

所组成的集合，其中 \(t_i\) 取遍一切实数值，而 \(n_j\) 取遍一切整数值。于是 \((\boldsymbol{u}|\boldsymbol{x})\) 对所有这些点 \(\boldsymbol{x}\) 都是整数，当且仅当 \((\boldsymbol{u}|\boldsymbol{a}_i) = 0\) （对 \(1 \leqslant i \leqslant p\) ）且 \((\boldsymbol{u}|\boldsymbol{a}_i)\) 是整数（对 \(p + 1 \leqslant i \leqslant p + q\) ）。用 \((\boldsymbol{a}_i')_{1 \leqslant i \leqslant n}\) 表示 \(\mathbf{R}^n\) 的满足 \((\boldsymbol{a}_i'|\boldsymbol{a}_j) = 0\) （对 \(i \neq j\) ）的基

\[
\left(\boldsymbol {a} _ {i} ^ {\prime} \mid \boldsymbol {a} _ {i}\right) = \mathrm{I} \quad \text { for } \quad \mathrm{I} \leqslant i \leqslant n
\]

{[}即与 \((a_{i})\) “对偶”的基{]}；若置 \(u = \sum_{i=1}^{n} u_{i} a_{i}^{\prime}\) ，则显然点 \(u \in G^{*}\) 由下列条件刻画： \(u_{i} = 0\) （对 \(i \leqslant i \leqslant p\) ），且 \(u_{i} \in Z\) （对

\[
p + \mathrm{i} \leqslant i \leqslant p + q;
\]

因而 \(G^{*}\) 是一个向量子空间 W 与一个离散子群的直和，前者的基由那些 \(a_{i}^{\prime}\) （满足 \(p + q + i \leqslant i \leqslant n\) ）组成，后者由那些 \(a_{i}^{\prime}\) （满足 \(p + i \leqslant i \leqslant p + q\) ）生成。换言之：

命题 5. 对 \(\mathbf{R}^n\) 的每个闭子群 G ，

\[
r (\mathbf {G} ^ {*}) = n - d (\mathbf {G}) \quad a n d \quad d (\mathbf {G} ^ {*}) = n - r (\mathbf {G}).
\]

让我们把同样的推理应用于 \(\mathbf{G}^*\) 。注意到与 \((\pmb{a}_i')\) 对偶的基是 \((\pmb{a}_i)\) ，我们就看出：

命题 6. 对每个子群 \(\mathbf{G}\) （ \(\mathbf{R}^n\) 的），我们有 \((\mathbf{G}^*)^* = \overline{\mathbf{G}}\) 。

推论. 一点 x 位于 \(R^{n}\) 的子群 G 的闭包中，当且仅当 \((u|x)\) 对每个 \(u \in R^{n}\) 都是整数，其中 \((u|y)\) 对一切 \(y \in G\) 均为整数。

让我们把这个关于位于子群 \(\mathbf{G}\) 的闭包中的点的刻画，应用于子群 \(\mathbf{G}\) 的情形——它由 \(n\) 个向量 \(\boldsymbol{e}_j\) （ \(\mathrm{i} \leqslant j \leqslant n\) ，属于典范基）以及任意 \(m\) 个点 \(\boldsymbol{a}_i\) （ \(\mathrm{i} \leqslant i \leqslant m\) ，属于 \(\mathbf{R}^n\) ）生成。说 \((\boldsymbol{u}|\boldsymbol{e}_j)\) 对 \(\mathrm{i} \leqslant j \leqslant n\) 是整数，就意味着 \(n\) 个坐标（ \(\boldsymbol{u}\) 的）都是整数；因此：

命题 7（克罗内克）. 设 \(\boldsymbol{a}_i = (a_{ji})\) （ \(i \leqslant i \leqslant m\) ， \(i \leqslant j \leqslant n\) ）是 \(m\) 个点（ \(\mathbf{R}^n\) 的）， \(\boldsymbol{b} = (b_j)\) （ \(i \leqslant j \leqslant n\) ）是 \(\mathbf{R}^n\) 的一个点。为了对每个 \(\varepsilon > 0\) ，存在 \(m\) 个整数 \(g_i\) （ \(i \leqslant i \leqslant m\) ）与 \(n\) 个整数 \(p_j\) （ \(i \leqslant j \leqslant n\) ），使得

\[
\left| q _ {1} a _ {1 j} + q _ {2} a _ {2 j} + \dots + q _ {m} a _ {m j} - p _ {j} - b _ {j} \right| \leqslant \varepsilon \quad (\mathrm{I} \leqslant j \leqslant n)
\]

其必要充分条件是：对每个有限序列 \((r_j)\) （ \(1 \leqslant j \leqslant n\) ，由 \(n\) 个整数组成），只要 \(m\) 个数 \(\sum_{j=1}^{n} a_{ij} r_j\) （ \(1 \leqslant i \leqslant m\) ）都是整数，数 \(\sum_{j=1}^{n} b_j r_j\) 就也是整数。

推论 1. 为了对每个 \(\boldsymbol{x} = (x_{j}) (\mathrm{I} \leqslant j \leqslant n)\) 与每个 \(\varepsilon > 0\) ，存在 m 个整数 \(q_{i} (\mathrm{I} \leqslant i \leqslant m)\) 与 n 个整数

\(p_j\) (1 \(\leqslant j\leqslant n\) )，使得

\[
\left| q _ {1} a _ {1 j} + q _ {2} a _ {2 j} + \dots + q _ {m} a _ {m j} - p _ {j} - x _ {j} \right| \leqslant \varepsilon \quad (\mathrm{I} \leqslant j \leqslant n)
\]

其必要充分条件是：不存在有限序列 \((r_j)\) （由 \(n\) 个整数组成、不全为零），使得 \(m\) 个数 \(\sum_{j=1}^{n} a_{ij} r_j\) 中的每一个都是整数。因为若 \(G\) 在 \(\mathbf{R}^n\) 中稠密，即若 \(\overline{\mathbf{G}} = \mathbf{R}^n\) ，则 \(\mathbf{G}^* = \{o\}\) ，反之亦然。

特别地 \((m = \mathrm{i})\) ：

推论 2. 设 \(\theta_{1},\theta_{2},\ldots ,\theta_{n}\) 是 \(n\) 个实数。为了对给定的任意 \(n\) 个实数 \(x_{1},x_{2},\ldots ,x_{n}\) 以及一个实数 \(\varepsilon >0\) ，存在一个整数 \(q\) 与 \(n\) 个整数 \(p_j\) ，使得

\[
\left| q \theta_ {j} - p _ {j} - x _ {j} \right| \leqslant \varepsilon \quad (\mathrm{I} \leqslant j \leqslant n)
\]

其必要充分条件是：不存在形如 \(\sum_{j=1}^{n} r_j \theta_j = h\) 的关系，其中 \(r_j\) 是 \(n\) 个不全为零的整数，而 \(h\) 是整数{[}这特别蕴涵诸 \(\theta_j\) 以及诸比 \(\theta_j / \theta_k\) （ \(j \neq k\) ）都必须是无理数{]}。

我们可以把这个结果解释如下：对每个整数 \(q \in Z\) ，设 \(x_{q}\) 是以 \(q\theta_{j} - [q\theta_{j}]\) （ \(i \leqslant j \leqslant n\) ）为坐标的点；于是推论 2 给出了使点 \(x_{q}\) 所成的集合在立方体（即区间 {[}0, i{]} ——在 \(R^{n}\) 的诸因子中——的乘积）中稠密的一个必要充分条件。

命题 8. 若 \(G_{1}\) 、 \(G_{2}\) 是 \(R^{n}\) 的任意闭子群，则我们有

\[
(\mathrm{G} _ {1} + \mathrm{G} _ {2}) ^ {*} = \mathrm{G} _ {1} ^ {*} \cap \mathrm{G} _ {2} ^ {*} \quad a n d \quad (\mathrm{G} _ {1} \cap \mathrm{G} _ {2}) ^ {*} = \mathrm{G} _ {1} ^ {*} + \mathrm{G} _ {2} ^ {*}.
\]

实数 \((\pmb{u}|\pmb{x} + \pmb{y})\) 对一切 \(\pmb{x} \in \mathbf{G}_1\) 与一切 \(\pmb{y} \in \mathbf{G}_2\) 都是整数，当且仅当 \((\pmb{u}|\pmb{x})\) 对一切 \(\pmb{x} \in \mathbf{G}_1\) 是整数且 \((\pmb{u}|\pmb{y})\) 对一切 \(\pmb{y} \in \mathbf{G}_2\) 是整数，因为 \((\pmb{u}|\pmb{x} + \pmb{y}) = (\pmb{u}|\pmb{x}) + (\pmb{u}|\pmb{y})\) ；因此

\[
(\mathrm{G} _ {1} + \mathrm{G} _ {2}) ^ {*} = \mathrm{G} _ {1} ^ {*} \cap \mathrm{G} _ {2} ^ {*}
\]

对任意两个子群 \(G_{1}\) 、 \(G_{2}\) （ \(R^{n}\) 的）成立。现在若假设 \(G_{1}\) 与 \(G_{2}\) 是闭的，则由命题 6 有 \((\mathbf{G}_{1}^{*} + \mathbf{G}_{2}^{*})^{*} = \mathbf{G}_{1} \cap \mathbf{G}_{2}\) ；因此，取相伴子群并再次应用命题 6，就得到 \((\mathbf{G}_{1} \cap \mathbf{G}_{2})^{*} = \overline{\mathbf{G}_{1}^{*} + \mathbf{G}_{2}^{*}}\) 。

注. 设 \(\mathbf{G}_1, \mathbf{G}_2\) 是 \(\mathbf{R}^n\) 中（第 1 小节意义下）的两个格，且 \(\mathbf{G}_2 \subset \mathbf{G}_1\) ；则（命题 5） \(\mathbf{G}_1^*\) 与 \(\mathbf{G}_2^*\) 是 \(\mathbf{R}^n\) 中满足 \(\mathbf{G}_1^* \subset \mathbf{G}_2^*\) 的格。我们在第 1 小节中已经看到，存在整数 \(m > 0\) 使 \(m\mathbf{G}_1 \subset \mathbf{G}_2\) ；

因而对 \(x \in G_1\) 与 \(u \in G_2^*\) ，我们有 \(m(u|x) \in Z\) ，从而 \((u|x) \in Q\) 。若 \(x \in G_2\) 且 \(u \in G_2^*\) ，或者 \(x \in G_1\) 且 \(u \in G_1^*\) ，则按定义有 \((u|x) \in Z\) 。于是，过渡到商后， \(Z\) -双线性映射 \((x, u) \to (u|x)\) （ \(G_1 \times G_2^*\) 到 \(Q\) 中的）定义了 \(Z\) -双线性映射 \(B\) （ \((G_1/G_2) \times (G_2^*/G_1^*)\) 到 \(Q/Z\) 中的）。此外，显然若 \(\bar{x}_0 \in G_1/G_2\) （相应地， \(\bar{u}_0 \in G_2^*/G_1^*\) ）使得对每个 \(\bar{u} \in G_2^*/G_1^*\) （相应地，对每个 \(\bar{x} \in G_1/G_2\) ）我们有

\[
\mathrm{B} (\overline {{{\boldsymbol {x}}}} _ {0}, \boldsymbol {u}) = \mathrm{o} \quad [ \text { resp. } \mathrm{B} (\overline {{{\boldsymbol {x}}}}, \overline {{{\boldsymbol {u}}}} _ {0}) = \mathrm{o} ],
\]

则必有 \(\overline{\mathbf{x}}_0 = \mathbf{o}\) （相应地， \(\overline{\mathbf{u}}_0 = \mathbf{o}\) ）。由此得出，存在一个 \(\mathbf{Z}\) -线性双射 \(h\) （ \(\mathrm{G}_2^* / \mathrm{G}_1^*\) 到 \(\mathrm{G}_1 / \mathrm{G}_2\) 的对偶上的），使得 \(\langle \overline{\mathbf{x}}, h(\overline{\mathbf{u}}) \rangle = \mathrm{B}(\overline{\mathbf{x}}, \overline{\mathbf{u}})\) （对 \(\overline{\mathbf{x}} \in \mathrm{G}_1 / \mathrm{G}_2\) 与 \(\overline{\mathbf{u}} \in \mathrm{G}_2^* / \mathrm{G}_1^*\) ）；特别地，有限群 \(\mathrm{G}_1 / \mathrm{G}_2\) 与 \(\mathrm{G}_2^* / \mathrm{G}_1^*\) 是同构的。

